% ===================================================================
\section{Apertura y planteamiento}

\diapo{0}{Portada}{30 s}
\begin{decir}
Buenos días. Presento PULSO, sigla de Pagos, Umbrales, Liquidez y Sincronización
Operativa. Es el proyecto del Semillero de Modelamiento de Sistemas Complejos y Física
Computacional, dirigido por el profesor Jhon Hernán Díaz Forero.

El proyecto estudia cómo los problemas de liquidez se propagan en una red de pagos
digitales. También estudia si el ritmo con que llegan los pagos, y no solo su monto,
permite anticipar el impago de un crédito.
\end{decir}
\begin{terminos}
\begin{description}
\item[Liquidez.] Disponibilidad de dinero en el momento en que hay que pagar. Se
distingue de la \emph{solvencia}, que compara el total de activos con el total de
deudas. Una entidad solvente puede quedarse sin liquidez si el dinero le llega tarde.
\item[Umbral.] Nivel de déficit que una entidad puede absorber antes de dejar de pagar.
En el modelo representa el colchón de liquidez.
\item[Sincronización operativa.] Grado en que los ciclos de pago de las entidades
coinciden en el tiempo.
\end{description}
\end{terminos}

\diapo{1}{Contenido}{20 s}
\begin{decir}
La exposición sigue el orden en que se construyó el proyecto. Primero planteo el
problema y la pregunta. Luego presento el modelo y las cinco tareas de la propuesta del
semillero, que van de la red de pagos a la predicción con aprendizaje automático.
Después presento la tarea 6, que aplica todo al riesgo crediticio de un banco. Termino
con los resultados, la implementación y las conclusiones.
\end{decir}

\diapo{2}{Contexto}{1 min}
\begin{decir}
Hoy casi todos los pagos de una economía pasan por sistemas digitales: salarios,
facturas entre empresas, arriendos y cuotas de crédito. En Colombia el Banco de la
República puso en marcha Bre-B, un sistema de pagos inmediatos que conecta a todas las
entidades financieras.

Cada persona o empresa paga con el dinero que recibe de otras. Por eso los pagos forman
una red. En la red, cada nodo es un agente y cada flecha es un pago de un agente a otro.
La flecha tiene un peso, que es el monto, y una dirección, que va de quien paga a quien
recibe.

En el diagrama, A tiene un problema y deja de pagar. B y C pierden parte de su ingreso.
Si no tienen reservas suficientes, tampoco pueden pagar a D. El problema de A llega a D
aunque D no tenga ninguna relación directa con A.

La inmediatez es una ventaja para los usuarios, pero también hace que un problema de
liquidez se transmita más rápido. Este proyecto modela esa transmisión.
\end{decir}
\begin{terminos}
\begin{description}
\item[Sistema de pagos.] Infraestructura que transfiere dinero entre cuentas de
distintas entidades y liquida las obligaciones.
\item[Bre-B.] Sistema de pagos inmediatos de bajo valor del Banco de la República de
Colombia, interoperable entre bancos y billeteras. En el proyecto aparece solo como
escenario de referencia. Los parámetros no están calibrados con sus datos.
\item[Red dirigida y ponderada.] Conjunto de nodos unidos por flechas con dirección y
peso. Se representa con una matriz $W$, donde $W_{ij}$ es el monto que $j$ paga a $i$.
\item[Impago.] Incumplimiento de una obligación de pago.
\end{description}
\end{terminos}
\begin{preguntas}
\pr{¿El proyecto usa datos de Bre-B?}{No. Todos los datos son simulados. Bre-B sirve
como motivación y como escenario ilustrativo de la aplicación (una red núcleo-periferia
de 245 entidades), con parámetros que no están calibrados.}
\end{preguntas}

\diapo{3}{Ciclos de pago y desfase}{1 min}
\begin{decir}
Los pagos no llegan de manera continua: siguen ciclos. La nómina llega cada quince o
treinta días, las facturas tienen plazos y las cuotas tienen fechas fijas.

Lo importante es la coordinación. Si el salario llega el día 30 y la cuota vence el día
5, el cliente tiene el dinero a tiempo. Si la nómina se corre al día 10, el monto es el
mismo pero el cliente necesita un colchón para cubrir esos días. A ese desfase lo llamamos
desincronización.

La gráfica de la izquierda muestra el mismo cliente con el mismo ruido, con y sin costo
por desfase. Con desfase, el déficit es mayor en todo momento. La gráfica de la derecha
muestra un empleado cuyo ritmo natural es 8~\% más rápido que el de su nómina. Con un
acoplamiento fuerte, el empleado ajusta su ritmo y se sincroniza; con un acoplamiento
débil, se desfasa cada vez más.

La idea central del proyecto es que el ingreso desfasado genera estrés financiero aunque
el monto no cambie.
\end{decir}
\begin{terminos}
\begin{description}
\item[Ciclo de pagos.] Periodo que se repite entre dos pagos del mismo tipo, por ejemplo
un mes de nómina.
\item[Fase.] Posición dentro del ciclo, medida como un ángulo entre $0$ y $2\pi$. La
fase $0$ es el inicio del ciclo y $2\pi$ es su final.
\item[Desfase.] Diferencia entre la fase de quien paga y la de quien recibe.
\item[Déficit.] Faltante de liquidez en un momento dado. En el modelo es la variable
$x_i$.
\item[Colchón.] Reserva que permite cubrir un déficit temporal sin dejar de pagar.
\item[Acoplamiento.] Intensidad con que un agente ajusta su ritmo al de quienes le
pagan.
\end{description}
\end{terminos}

\diapo{4}{Limitación del score clásico y pregunta de investigación}{1 min}
\begin{decir}
Los bancos deciden a quién prestar con un score crediticio. El score es un modelo
estadístico que estima la probabilidad de impago con variables del propio cliente:
ingreso, endeudamiento, historial de mora, antigüedad y tipo de empleo.

El score evalúa a cada cliente de forma aislada. No ve tres cosas: si quienes le pagan
al cliente están en problemas, si todo su ingreso depende de un solo pagador, y si sus
ingresos llegan desfasados de sus obligaciones.

Esa información existe en los registros de transacciones del banco, pero el score
clásico no la usa.

La pregunta del proyecto es esta: la sincronización de los ciclos de pago, estimada
únicamente con las transacciones, ¿aporta información sobre el impago que no esté ya en
el score clásico ni en la estructura estática de la red? La comparación con la red
estática es la parte exigente: ya existen trabajos que usan la red de un cliente para
mejorar el score. Lo nuevo es la dinámica.
\end{decir}
\begin{terminos}
\begin{description}
\item[Score crediticio.] Modelo estadístico que asigna a cada cliente una probabilidad
de impago, normalmente con regresión logística sobre variables individuales.
\item[PD.] Probabilidad de incumplimiento (\emph{probability of default}): probabilidad
de que el cliente deje de pagar su crédito en un horizonte, aquí doce meses.
\item[Red estática.] Variables que describen quién le paga a quién y cuánto, sin
considerar el momento de los pagos: concentración del ingreso, número de pagadores,
centralidad.
\item[Red dinámica.] Variables que describen el ritmo: sincronización, desfase,
regularidad de los ciclos.
\end{description}
\end{terminos}
\begin{preguntas}
\pr{¿No es esto lo que ya hacen los bancos con el historial crediticio?}{El historial
describe al cliente por sí solo. Las variables de red estática describen sus relaciones
y ya aparecen en la literatura (por ejemplo, Óskarsdóttir y Bravo, 2021). La
sincronización de los ciclos de pago no aparece en esos trabajos, y por eso el
experimento la compara contra ambos niveles.}
\end{preguntas}

\diapo{5}{Hipótesis y objetivos}{1 min}
\begin{decir}
La hipótesis es falsable, es decir, puede resultar falsa. Si el ingreso desfasado
genera déficit, las variables de sincronización deben mejorar la predicción del impago.
Si ese mecanismo no existe, no deben mejorarla. Por eso el experimento incluye un
control nulo: un mundo idéntico donde la red no genera déficit. Si las variables
mejoraran la predicción también ahí, el resultado sería un artefacto del método.

El objetivo general es construir, paso a paso, un modelo físico-matemático de la red de
pagos y aplicarlo al riesgo crediticio. La tabla muestra las seis tareas. Las cinco
primeras vienen de la propuesta del semillero: déficit en la red, difusión,
sincronización, impagos en cascada y aprendizaje automático. La sexta es la aplicación a
la cartera de un banco.
\end{decir}
\begin{terminos}
\begin{description}
\item[Hipótesis falsable.] Afirmación que predice un resultado observable que podría no
ocurrir; si no ocurre, la hipótesis se rechaza.
\item[Control nulo.] Experimento idéntico al principal en el que se apaga el mecanismo
estudiado. Mide cuánto ``efecto'' produce el método por sí solo.
\item[Kuramoto.] Modelo de osciladores acoplados propuesto por Yoshiki Kuramoto (1975)
para estudiar la sincronización. Se explica en la diapositiva 9.
\end{description}
\end{terminos}

% ===================================================================
\section{Modelo unificado}

\diapo{6}{Un modelo, seis tareas}{2 min}
\begin{decir}
Esta es la ecuación central del proyecto. Cada entidad $i$ tiene dos variables: su
déficit $x_i$ y la fase de su ciclo de pagos $\theta_i$.

La primera ecuación dice cómo cambia el déficit en el tiempo. Tiene cinco términos.
Primero, $-\alpha x_i$: cada entidad corrige su déficit por sí misma, por ejemplo con
ahorro o crédito de corto plazo. Segundo, $\beta\sum_j W_{ij}x_j$: si quienes le pagan
están en déficit, le pagan menos y su déficit aumenta; este es el contagio. Tercero,
$-D_0(Lx)_i$: la redistribución de liquidez entre entidades, que traslada déficit sin
crearlo. Cuarto, $\gamma k_i(1-r_i)$: el ingreso que llega desfasado se vuelve déficit.
Quinto, $\eta_i(t)$: los choques externos.

La segunda ecuación describe el ritmo de pagos. Cada entidad tiene su ritmo natural
$\omega_i$ y lo ajusta según el ritmo de quienes le pagan, con intensidad $K$.

Finalmente, si el déficit supera el umbral $x_i^{(c)}$, la entidad entra en impago y deja
de pagar.

Las etiquetas T1 a T4 indican en qué tarea entra cada término. El criterio de diseño es
que, si se apagan los términos nuevos, el modelo coincide exactamente con el de la tarea
anterior. Eso se verificó con pruebas automáticas. Así, cada resultado se puede
comparar con el anterior.
\end{decir}
\begin{terminos}
\begin{description}
\item[$x_i(t)$.] Déficit de liquidez de la entidad $i$ en el instante $t$. En las
tareas 1–5 se mide en ``pagos medios por ciclo''.
\item[$t$.] Tiempo. En las tareas 1–5, una unidad es un ciclo de pagos (un día hábil);
en la tarea 6, un mes.
\item[$dx_i/dt$.] Tasa de cambio del déficit: cuánto crece o decrece por unidad de
tiempo.
\item[$\alpha$.] Tasa de autocorrección [1/ciclo]. Con solo este término el déficit
decae como $e^{-\alpha t}$; $1/\alpha$ es el tiempo típico de corrección.
\item[$\beta$.] Intensidad del contagio [1/(ciclo·pago)]: déficit que genera en $i$
cada unidad de déficit de un pagador, por unidad de pago.
\item[$W_{ij}$.] Pago medio de $j$ hacia $i$ por ciclo. Las filas de $W$ son ingresos
y las columnas son egresos.
\item[$\sum_j W_{ij}x_j$.] Suma del déficit de todos los pagadores de $i$, ponderada
por cuánto le paga cada uno.
\item[$D_0$.] Coeficiente de redistribución [1/(ciclo·pago)].
\item[$L$.] Laplaciano de flujo, $L=D_\text{out}-W$ (diapositiva 8).
\item[$(Lx)_i$.] Componente $i$ del vector $Lx$.
\item[$\gamma$.] Peso del desfase en el déficit [1/ciclo].
\item[$k_i$.] Ingreso por la red: $k_i=\sum_jW_{ij}$.
\item[$r_i$.] Coherencia local: entre $-1$ (ingreso en contrafase) y $1$ (ingreso en
fase).
\item[$\eta_i(t)$.] Forzamiento externo: choques de liquidez en las tareas 1–5 y ruido
aleatorio en la tarea 6.
\item[$\theta_i$.] Fase del ciclo de pagos de $i$, en radianes.
\item[$\omega_i$.] Frecuencia natural [rad/ciclo]: ritmo de $i$ si no tuviera pagadores.
\item[$K$.] Intensidad del acoplamiento de fase [rad/ciclo].
\item[$\langle k\rangle$.] Ingreso medio por entidad. Normaliza $K$ (diapositiva 9).
\item[$\sin(\theta_j-\theta_i)$.] Positivo si el pagador $j$ va adelantado: acelera a
$i$. Negativo si va atrasado: lo frena.
\item[$x_i^{(c)}$.] Umbral de impago (colchón).
\item[Criterio de anidamiento.] Con $D_0=\gamma=K=0$ y sin umbral se recupera la
tarea 1; con $\gamma=K=0$, la tarea 2; y así sucesivamente.
\end{description}
\end{terminos}
\begin{preguntas}
\pr{¿Por qué un modelo lineal en $x$?}{Es la aproximación más simple que conserva el
contagio y permite resultados analíticos: el umbral de estabilidad $R=1$, la
conservación bajo difusión y la tasa de recuperación. La no linealidad entra por el
umbral y por las fases. Es una limitación declarada.}
\pr{¿El déficit puede afectar el ritmo de pagos?}{En esta versión no: el acople va del
ritmo al déficit. La retroalimentación inversa aparece en el trabajo futuro.}
\end{preguntas}
